\documentclass[11pt]{article}
\usepackage[margin=1in]{geometry}
\usepackage[T1]{fontenc}
\usepackage{lmodern}
\usepackage{amsmath,amssymb,amsthm,mathtools}
\usepackage{booktabs,array,longtable,xcolor,listings,xurl}
\usepackage[hidelinks]{hyperref}
\hypersetup{pdftitle={An Absorbing Set Without Gallager-B Trapping},pdfauthor={}}
\newtheorem{theorem}{Theorem}
\newtheorem{lemma}[theorem]{Lemma}
\newtheorem{proposition}[theorem]{Proposition}
\newcommand{\F}{\mathbb F_2}
\newcommand{\supp}{\operatorname{supp}}
\newcommand{\row}{\operatorname{row}}
\newcommand{\maj}{\operatorname{maj}_0}
\setlength{\emergencystretch}{3em}
\lstset{basicstyle=\scriptsize\ttfamily,columns=fullflexible,keepspaces=true,
 breaklines=true,breakatwhitespace=false,showstringspaces=false,frame=lines}
\title{An Absorbing Set Without Gallager-B Trapping}
\author{}
\date{October 26, 2026}
\begin{document}
\maketitle
\begin{abstract}
We give a seven-variable $(7,0)$-absorbing set for which every binary error
support is accepted by the zero-initialized syndrome-based Gallager-B decoder
at initialization or after one iteration. Both the hard-decision syndrome and
the parity of outgoing extrinsic messages match the input syndrome. Thus no
subset produces the persistent syndrome mismatch or nonconvergence in the
trapping-set definition of Morris, Pllaha and Kelley. The matrix consists of
all seven nonzero simplex checks: majority voting recovers the single-column
representative of each Hamming syndrome. To distinguish syndrome acceptance
from quantum recovery, we also give an explicit fourteen-qubit CSS pair in
which this absorbing component corrects every error in the decoded binary
sector modulo the entire opposite stabilizer span. A nontrivial logical
operator remains on the other component. Complete finite proofs and the
stopped-execution argument are checked in Lean, with an independent
literal-matrix verifier.
\end{abstract}

\section{Question and answer}
Morris, Pllaha and Kelley ask in Section VI of
\emph{Absorbing Sets in Quantum LDPC Codes}, arXiv:2307.14532v2,
May 10, 2024, whether every absorbing set is a trapping set for their
syndrome-based Gallager-B decoder \cite{mpk}.
Their Definition 5 requires each selected variable to have strictly more
neighbors of even induced degree than of odd induced degree.
Definition 3 calls a nonempty set trapping if some error supported on it
produces a neighboring check that is not eventually correct or a selected
variable that does not eventually converge.

\begin{theorem}\label{thm:main}
There is a connected seven-variable $(7,0)$-absorbing Tanner graph such
that every error supported on its variables has an accepted, constant
terminal output after at most one full Gallager-B iteration. Consequently
the variable set is not trapping in the sense of Definition 3 of
\cite{mpk}. Moreover it occurs as a component of an explicit CSS code with
a nontrivial logical operator, in which every error in this binary sector
supported on the component is corrected modulo stabilizers.
\end{theorem}
The first assertion is proved in Section~\ref{sec:decoding}; the CSS
assertion is proved in Section~\ref{sec:css}. The precise source-to-model
comparison is in Section~\ref{sec:meaning}. Appendix~\ref{app:lean} contains
central Lean definitions, complete proof bodies and direct source links.

The example uses redundant checks, check degree four, variable degree four
and Tanner four-cycles. None is forbidden in the cited general question.
The selected component is connected; the fourteen-qubit ambient graph has
two components. We do not assert a connected-ambient, girth-constrained or
column-weight-three counterexample. We also do not assert correction of
both Pauli sectors on the selected component, or of arbitrary errors on
the other component.

\section{The actual matrix}
Index checks and variables by the nonzero vectors of $\F^3$, in the
integer order $1,\ldots,7$ with least significant binary digit first.
Let $H_{a,v}=a\cdot v$. Explicitly,
\begin{equation}\label{eq:H}
H=\begin{pmatrix}
1&0&1&0&1&0&1\\
0&1&1&0&0&1&1\\
1&1&0&0&1&1&0\\
0&0&0&1&1&1&1\\
1&0&1&1&0&1&0\\
0&1&1&1&1&0&0\\
1&1&0&1&0&0&1
\end{pmatrix}.
\end{equation}
Write $A$ for all seven variables. Every row and column of $H$ has weight
four, so all neighboring checks have even induced degree and every variable
has four even and no odd neighbors. Hence $A$ is $(7,0)$-absorbing.
Any pair of distinct variables has two common check neighbors: two independent
linear equations $a\cdot u=a\cdot v=1$ have two solutions in $\F^3$.
In particular the Tanner graph is connected and has four-cycles.
The finite statements \texttt{degrees}, \texttt{full\_set\_absorbing} and
\texttt{first\_block\_connected} check the actual matrix, not an abstract
replacement graph.

\section{One iteration accepts every supported error}\label{sec:decoding}
All arithmetic in check messages is in $\F$.
For an arbitrary input error $e\in\F^7$ put
\[
 p=\sum_{v\ne0}e_vv\in\F^3,\qquad s=He^{\mathsf T}.
\]
Then $s_a=a\cdot p$. Zero initial variable-to-check messages imply that
check $a$ sends $s_a$ to each of its neighbors. A hard decision is the
strict majority of all incoming check messages, with ties resolved to zero.
The outgoing message toward $a$ takes the same strict majority after
excluding $a$.

\begin{lemma}\label{lem:votes}
For $v\ne0$, the number of neighboring checks sending one initially is
zero if $p=0$, four if $p=v$, and two otherwise. Thus the first hard output is
\[
 \hat e=0\quad(p=0),\qquad \hat e=\mathbf e_p\quad(p\ne0),
\]
where $\mathbf e_p$ is the singleton indicator at column $p$.
In every case $H\hat e^{\mathsf T}=s$.
\end{lemma}
\begin{proof}
The vote count is the number of $a\in\F^3\setminus\{0\}$ satisfying
$a\cdot v=a\cdot p=1$. If $p=0$ there are none. If $p=v$ there are four.
Otherwise the two nonzero, distinct vectors $v,p$ are independent and the
two independent equations have two solutions. The strict majority of four
bits is one only in the second case. The asserted syndrome follows from
$H\mathbf e_p^{\mathsf T}=(a\cdot p)_a$.
All counts and the displayed output formula have matching kernel-checked
finite Lean proofs, including all $128$ inputs.
\end{proof}

The conclusion also survives the alternate syndrome calculation from outgoing
extrinsic variable messages, which is relevant to some computations in
\cite{mpk}.
\begin{lemma}\label{lem:edges}
Let $x'_{v\to a}$ be the first outgoing extrinsic message. Then
$\sum_{v:a\cdot v=1}x'_{v\to a}=s_a$ for every check $a$.
\end{lemma}
\begin{proof}
For $p=0$ all messages vanish. For $p\ne0$, variable $p$ sends one on all
four incident edges. At $v\ne p$ there were two ones and two zeros before
excluding the recipient check. Excluding a zero gives two ones among three
messages; excluding a one gives one. Hence $x'_{v\to a}=1-s_a$ for these
variables. A check with $s_a=1$ is adjacent to $p$ and receives exactly one
one. A check with $s_a=0$ is not adjacent to $p$ and receives four ones,
whose parity is zero. This is exactly $s_a$. The actual update equations
are checked by \texttt{first\_edge\_matches}.
\end{proof}

\begin{proof}[Proof of the dynamical assertion of Theorem~\ref{thm:main}]
The zero syndrome is accepted at initialization. Every nonzero syndrome
is accepted after the first iteration by Lemma~\ref{lem:votes};
Lemma~\ref{lem:edges} gives the same stopping decision for the alternate
edge-parity convention. Extend a halted decoder by repeating its terminal
output. All neighboring checks are then correct, and every variable is
constant, for all times $t\ge1$. This holds for every error $e$, not just
the error with support $A$. There can therefore be no failure-inducing
subset in Definition 3. The Lean theorem
\texttt{absorbing\_not\_trapping} proves precisely these eventual properties.
\end{proof}
The extension by a constant terminal output represents the stated stopping
algorithm. It does not claim that artificially continuing message updates
after a successful stop would preserve every intermediate message.

\section{An actual CSS counterexample with successful recovery}\label{sec:css}
Syndrome acceptance alone does not guarantee quantum recovery. For example,
using $H$ in both sectors of the familiar Steane construction \cite{steane},
errors on columns $1$ and $2$ produce the singleton output on column $3$.
The residual supported on $\{1,2,3\}$ has zero syndrome but is not in
$\row H$. Thus it is a logical miscorrection. The source's dynamical
Definition 3 and the broader notion of quantum decoding failure must be
distinguished. We now eliminate this issue in a separate explicit CSS pair.

Let $G$ consist of rows $1,2,4$ of $H$, and put
\[
 K=\begin{pmatrix}G\\\mathbf1\end{pmatrix},\qquad
 H_X=\begin{pmatrix}H&0\\0&H\end{pmatrix},\qquad
 H_Z=\begin{pmatrix}K&0\\0&G\end{pmatrix}.
\]
These are binary matrices of sizes $14\times14$ and $7\times14$.
The redundant rows of $H_X$ are retained as the actual decoding checks.
For $Z$-errors, the decoder uses $H_X$ and successful recovery means that
the residual belongs to the \emph{entire} row span of $H_Z$.

\begin{lemma}\label{lem:kernel}
$\ker H=\row K$, and $H_XH_Z^{\mathsf T}=0$.
\end{lemma}
\begin{proof}
Every row of $H$ has even weight, and two distinct rows have overlap two.
Consequently $HG^{\mathsf T}=0$ and $H\mathbf1^{\mathsf T}=0$, so
$\row K\subseteq\ker H$ and the displayed CSS matrices commute.
For completeness, the reverse inclusion is checked on all $128$ binary
words against all $16$ combinations of the four displayed generators.
The Lean theorem \texttt{kernel\_equals\_four\_generator\_span} is this
complete finite proof, not a search restricted by weight. The independent
literal-matrix verifier obtains exactly the same kernel of $16$ words.
\end{proof}

\begin{proposition}\label{prop:css}
Let $A$ be the first seven coordinates in this CSS pair. Every binary
$Z$-error supported on $A$ is decoded correctly modulo the complete
$Z$-stabilizer span. There is nevertheless a nontrivial $Z$ logical
operator on the second block.
\end{proposition}
\begin{proof}
There are no check edges between the two blocks. On input $(e,0)$ the
second block remains zero, while the first returns $(\hat e,0)$ at
initialization or after one iteration. By Lemma~\ref{lem:votes},
$e+\hat e\in\ker H=\row K$. Hence
$(e+\hat e,0)\in\row H_Z$, so recovery is quantum-correct for all $128$
supported errors. This exact complete-span membership is also checked
directly by \texttt{first\_block\_quantum\_recovery}.

The word $(0,\mathbf1)$ commutes with every row of $H_X$, because all rows
of $H$ have even weight. It is not in $\row H_Z$: all second-block
components in that span belong to $\row G$, whose words have even weight,
whereas $\mathbf1$ has weight seven. This supplies a genuine nontrivial
logical operator. Both statements are checked on the actual matrices by
\texttt{logical\_commutes} and \texttt{logical\_not\_stabilizer}.
\end{proof}
This construction explicitly uses a stabilizer-state component for the
selected support, together with a component carrying logical information.
It is not a claim that arbitrary high-weight errors on encoded information
can be corrected. It shows why an unrestricted assertion about every
absorbing set needs hypotheses on both the Tanner representation and the
opposite stabilizer space.

\section{Comparison with the source and earlier ingredients}\label{sec:meaning}
The decoder here uses exactly zero initialization, extrinsic check XOR,
extrinsic variable majority, zero ties, and syndrome stopping from Section
II-B of \cite{mpk}. Every possible subset of $A$ is an input. There is no
iteration cutoff, sampled failure estimate, or assumption excluding
stabilizer errors. The first block has no exterior check neighbors, so the
source convention that exterior messages are correct is automatically met.
The checks and variables relevant to Definition 3 are all included.

The source proves the positive-$b$ absorbing-set result in its Theorem 4;
our example has $b=0$ and does not contradict that theorem. Its path and
cycle results concern other graphs; our matrix has degree four and is
neither a path nor a cycle. In particular the counterexample does not
imply that those cases, or useful high-girth code-design restrictions,
are settled by this example.

Hamming's construction \cite{hamming} already provides the syndrome-column
representation used here. The simplex checks, the Steane code and the CSS
stabilizer criterion are classical ingredients. The contribution here is
their explicit application to the blanket absorbing-set question, the
separation of dynamical trapping from logical miscorrection, and the
complete checked counterexample. The two-block CSS example also refutes
the stronger interpretation that some supported error must fail quantum
recovery in the selected sector. It does not replace the opposite
stabilizer span by the row span of the decoded matrix.

The directly inspected question is the May 10, 2024 arXiv version. A journal
version appeared March 3, 2026 (DOI 10.3934/amc.2026023); its full text was
access-restricted during the comparison. We therefore do not attribute
unchanged wording to the journal version. The complete search and
source-screening ledger is linked in Appendix~\ref{app:lean}.

\begin{table}[ht]
\centering\small
\begin{tabular}{p{.26\textwidth}p{.34\textwidth}p{.30\textwidth}}
\toprule Source & Inspected milestone & Interval to October 26, 2026\\
\midrule
Hamming & April 1950 publication & About 76.5 years (month precision)\\
Steane & January 29, 1996 preprint & 11,228 days\\
Morris--Pllaha--Kelley & July 26, 2023 first preprint & 1,188 days\\
Question in inspected version & May 10, 2024, v2, Section VI & 899 days\\
Journal metadata & March 3, 2026 online publication & 237 days\\
This paper & October 26, 2026 & 0 days\\
\bottomrule
\end{tabular}
\caption{Reference chronology. Ingredient age is not the duration for which
an open question was unresolved. Exact day differences are independently
regenerated in the artifact.}
\end{table}

\section{Conclusion}
The full variable set of the redundant seven-column simplex matrix is
absorbing, but none of its subsets traps the stated stopped Gallager-B
algorithm. The explicit CSS pair additionally makes every supported error
in the decoded sector quantum-correct. Thus an unrestricted implication
from absorbing to trapping is false. Whether useful additional restrictions
such as connected ambient CSS graphs, prescribed degree or girth conditions,
or requiring the selected support to carry logical information restore such
an implication is a separate question.

\begin{thebibliography}{9}
\bibitem{mpk} K. D. Morris, T. Pllaha and C. A. Kelley,
\emph{Absorbing Sets in Quantum LDPC Codes}, arXiv:2307.14532v2,
May 10, 2024, Sections II-B, III (Definitions 3 and 5), IV and VI.
\url{https://arxiv.org/html/2307.14532v2}.
Journal metadata: \emph{Advances in Mathematics of Communications} 23
(2026), 234--254, \url{https://doi.org/10.3934/amc.2026023}.
\bibitem{hamming} R. W. Hamming, \emph{Error detecting and error correcting
codes}, Bell System Technical Journal 29 (1950), 147--160.
\url{https://doi.org/10.1002/j.1538-7305.1950.tb00463.x}.
\bibitem{steane} A. Steane, \emph{Multiple particle interference and quantum
error correction}, arXiv:quant-ph/9601029, January 29, 1996;
Proceedings of the Royal Society A 452 (1996), 2551--2577.
\url{https://arxiv.org/abs/quant-ph/9601029}.
\end{thebibliography}
\appendix
\section{Complete checked proofs and source bridges}\label{app:lean}
The single research module imports only \texttt{Std}. The proof toolchain is
Lean 4.36.0-rc1, commit \texttt{c9afcddf514a6772bd4d649f7cfc04eff0b9f15e}.
Every finite statement below is proved by kernel reduction with
\texttt{decide}, not by a native decision oracle. The quantifier
\texttt{e : BitVec 7} covers every one of the $128$ subsets, including the
empty and full supports. The clean build and axiom reports are included
in the artifact. The permitted axioms are \texttt{propext},
\texttt{Classical.choice} and \texttt{Quot.sound}; no admitted proof is used.

\subsection{Public sources and the proof map}
The \href{https://github.com/alectionary/vibe-math/blob/master/formalization/gallager-b-absorbing/Counterexample.lean}{complete current Lean source}
and its
\href{https://github.com/alectionary/vibe-math/blob/2717b0ab2b59c2ccb71f0e369a9f21d80490c191/formalization/gallager-b-absorbing/Counterexample.lean}{immutable checked revision}
contain all research proofs, including the actual fourteen-coordinate bridge.
The
\href{https://github.com/alectionary/vibe-math/blob/master/formalization/gallager-b-absorbing/verification.json}{verification manifest},
\href{https://github.com/alectionary/vibe-math/blob/master/formalization/gallager-b-absorbing/prior-art.json}{source and search ledger}, and
\href{https://github.com/alectionary/vibe-math/blob/master/formalization/gallager-b-absorbing/verify.py}{independent literal-matrix verifier}
make the input domain, access limits and evidence distinctions explicit.

\small
\begin{longtable}{p{.27\textwidth}p{.61\textwidth}}
\toprule Manuscript proof & Matching Lean declarations (all in the linked module)\\
\midrule\endhead
Absorption and connectivity & \texttt{degrees}, \texttt{full\_set\_absorbing}, \texttt{first\_block\_connected}\\
Syndrome-column identity and Lemma~\ref{lem:votes} & \texttt{syndrome\_index\_bridge}, \texttt{simplex\_vote\_counts}, \texttt{first\_hard\_formula}, \texttt{first\_hard\_matches}\\
Extrinsic update, Lemma~\ref{lem:edges} & \texttt{first\_edge\_matches}, with the full \texttt{checkUpdate} and \texttt{variableUpdate} definitions\\
Theorem~\ref{thm:main}, eventual properties & \texttt{terminal\_matches}, \texttt{all\_eventually\_correct}, \texttt{all\_eventually\_converge}, \texttt{absorbing\_not\_trapping}\\
Steane miscorrection comparison & \texttt{steane\_logical\_miscorrection}\\
Lemma~\ref{lem:kernel} & \texttt{kernel\_equals\_four\_generator\_span}, \texttt{actual\_css\_commutation}\\
Proposition~\ref{prop:css}, complete span and nontrivial logical & \texttt{first\_block\_residual\_is\_stabilizer}, \texttt{first\_block\_quantum\_recovery}, \texttt{logical\_commutes}, \texttt{logical\_not\_stabilizer}\\
Actual source matrices and fourteen-coordinate update & \texttt{actual\_matrix\_first\_step\_bridge}, \texttt{actual\_matrix\_syndrome\_matches}, \texttt{terminal\_eq\_firstHard}, \texttt{actual\_matrix\_quantum\_recovery}\\
\bottomrule
\end{longtable}
\normalsize

\lstset{escapeinside={(*@}{@*)}}

\subsection{Matrix, subsets and exact update rules}
A seven-bit word is the indicator of a subset of the seven actual columns.
The matrix is the binary dot product on labels $1,\ldots,7$.
\texttt{pack\_bits} checks that packing and reading these column bits agree.
The two message functions exclude the recipient explicitly, and the
majority definition implements zero ties by a strict inequality.
Here are the complete definitions and finite source checks. Unicode symbols are rendered typographically; the linked Lean file is the checked source.
\lstinputlisting[firstline=1,lastline=81]{proof-listing.txt}

\subsection{From accepted output to absence of trapping}
The dependency chain is
\texttt{first\_hard\_matches} $\to$ \texttt{terminal\_matches} $\to$
the two eventual properties $\to$ \texttt{absorbing\_not\_trapping}.
The output at time zero is the initial zero guess. From time one onward
it is the accepted terminal value. Thus the following unbounded-in-time
proof is a structural argument, rather than a finite simulation horizon.
The complete central proof bodies are printed here.
\lstinputlisting[firstline=83,lastline=126]{proof-listing.txt}
Direct declaration links:
\href{https://github.com/alectionary/vibe-math/blob/2717b0ab2b59c2ccb71f0e369a9f21d80490c191/formalization/gallager-b-absorbing/Counterexample.lean#L101}{eventual check correctness},
\href{https://github.com/alectionary/vibe-math/blob/2717b0ab2b59c2ccb71f0e369a9f21d80490c191/formalization/gallager-b-absorbing/Counterexample.lean#L108}{eventual variable convergence}, and
\href{https://github.com/alectionary/vibe-math/blob/2717b0ab2b59c2ccb71f0e369a9f21d80490c191/formalization/gallager-b-absorbing/Counterexample.lean#L115}{nontrapping theorem}.

\subsection{Full CSS span and logical witness}
A pair of seven-bit words is exactly a binary word on the disjoint union
of two seven-coordinate blocks. The matrices below retain all fourteen
$X$ checks and all seven $Z$ generators. \texttt{combine} enumerates the
entire generator span with independent binary coefficients; it does not
check only individual stabilizer rows. The recovery statement covers every
input and witnesses membership of the residual in that whole span.
The nontrivial logical test uses the same matrix and span definitions.
\lstinputlisting[firstline=128,lastline=191]{proof-listing.txt}
Direct declaration links:
\href{https://github.com/alectionary/vibe-math/blob/2717b0ab2b59c2ccb71f0e369a9f21d80490c191/formalization/gallager-b-absorbing/Counterexample.lean#L151}{complete residual-span theorem} and
\href{https://github.com/alectionary/vibe-math/blob/2717b0ab2b59c2ccb71f0e369a9f21d80490c191/formalization/gallager-b-absorbing/Counterexample.lean#L165}{nontrivial logical theorem}.

\subsection{Analytic identities and the actual fourteen-coordinate decoder}
The finite formulas below supply the syndrome-column, vote-count and
kernel/span identities used in the written proofs. The final portion then
calculates the syndrome and all first-step votes directly from the actual
fourteen-coordinate rows. Its bridge theorem proves that this calculation
is the component calculation used above. This avoids assuming that an
abstract block argument automatically matches the source matrix.
\lstinputlisting[firstline=193,lastline=274]{proof-listing.txt}
Direct declaration link:
\href{https://github.com/alectionary/vibe-math/blob/2717b0ab2b59c2ccb71f0e369a9f21d80490c191/formalization/gallager-b-absorbing/Counterexample.lean#L259}{actual-matrix quantum recovery theorem}.

\subsection{Reproduction and evidence limits}
Run \texttt{python3 build.py --lean /path/to/lean} from the artifact
folder with the pinned toolchain, then \texttt{python3 verify.py --check}.
The Lean result is the proof; the independently regenerated JSON certificate
provides a second, literal-matrix reconstruction of all inputs.
No sampled experiment is used. The certificate also supplies independent
linear-algebra diagnostics; numerical diagnostics not stated as theorems
in the manuscript are not additional formal results.
The source ledger distinguishes the accessible arXiv source, the restricted
journal text, partial classical-source retrievals and historical questions
harvested without claiming they remain open today.

\end{document}
